\section{Interlude (Act~I, confirmed): adapting released guards --- a preregistered starting-type study}
\label{sec:adaptation}

Acts~I--II characterize what fine-tuning does to \emph{general} instruction checkpoints, and Act~III ranks
\emph{base} guards on regulated domains --- all on panels the researcher inspected while building the
method, so they are retrospective estimation, not confirmed findings (\Cref{sec:honesty}). This section is
the one \emph{analysis-preregistered} piece: its estimands, decision rules, non-inferiority margin, and
interpretation wording were fixed in a committed claim registry
(\code{artifacts/starting\_type\_adaptation\_v1/protocol/claim\_registry.json}) \emph{before} any score
existed --- git history puts the registry a day ahead of the scores --- so the verdicts below cannot be a
post-hoc reading of the data (no HARKing). Two limits on that label, stated here rather than in the
appendix. The registry declares itself \code{finalization\_status: dev\_nonfinal} and is \emph{not} bound
to a release lock (no lock exists for this study). And the study re-scores the \emph{same 3{,}308 rows} as
Acts~I--II, from the same frozen manifest --- so it is preregistered on the \emph{analysis}, not blind on
the \emph{data}; the uninspected-cohort half of that discipline remains future work
(\Cref{sec:roadmap}). It asks two questions that the earlier acts raise but cannot settle:
(RQ1) does the specialization tradeoff also hit models that are \emph{already} purpose-built safety
guards, and (RQ2) is KL-SFT --- which recovered transfer for general checkpoints (\Cref{sec:actI-klsft})
--- a \emph{free} improvement, i.e.\ does it retain transfer at no represented-source cost?

\subsection{Design: a $2\times3$ blocked grid over ten checkpoints}
\label{sec:adaptation-design}
\noindent A $2\times 3$ blocked grid: $\{$four general instruction checkpoints, six released
purpose-built guards$\}\times\{$unmodified $U$, $+$SFT, $+$KL-SFT$\}$, over \AdaNCheckpoints{}
checkpoints spanning \AdaNFamilies{} model families (gemma, granite, llama, mistral, qwen, smollm), five
seeds, primary $\beta{=}\AdaPrimaryBeta$. Every guard keeps its \emph{native} top-level verdict
interface (ShieldGemma \texttt{Yes/No}, Qwen3Guard \texttt{Safe/Controversial/Unsafe} --- its native
three-tier top-level verdict, reduced to the same binary decision margin as every other cell,
$z_{\text{Unsafe}}-z_{\text{Safe}}$, which is what the committed scores record (the
\texttt{Controversial} logit is not used) --- Llama-Guard
\texttt{safe/unsafe}, Granite \texttt{Yes/No}, WildGuard \texttt{yes/no}), validated byte-for-byte
against the real tokenizer at a decision-position fidelity check before scoring; all cells share the
frozen Paper~A binary manifest, LoRA recipe, and optimizer. We report macro-AP (mean over benchmark
sources) on the RAW logit margin, on represented (\texttt{id\_test}) and held-out (\texttt{transfer})
splits, as a within-checkpoint change vs.\ the \emph{same} unmodified checkpoint (the KL reference is
that checkpoint). Each statistic H is an equal-weight mean over the six \emph{model} families; its
one-sided $97.5\%$ lower bound comes from a 10{,}000-resample bootstrap that resamples near-duplicate
\emph{evaluation row families} (\code{family\_id}; 3{,}170 clusters over 3{,}308 rows) and training seeds while
holding model \emph{and benchmark-source} identity fixed, Bonferroni-split
across the two research questions (familywise $\alpha=0.05$). Because benchmark identity is held fixed,
these bounds carry evaluation-row and seed uncertainty only --- not between-benchmark or
between-model-family uncertainty.

\input{generated/tab_adaptation_gen}

\begin{figure}[htbp]\centering
\includegraphics[width=0.86\linewidth]{fig_adaptation_plane.pdf}
\caption{\textbf{The adaptation specialization plane.} Each checkpoint's macro-AP change vs.\ its
\emph{own} unmodified base, on represented (x) and held-out transfer (y) sources, under SFT ($\circ$)
and KL-SFT ($\triangle$); arrows run SFT\,$\to$\,KL-SFT. Nearly all checkpoints --- general (blue) and
released purpose-built guards (green) alike --- land in the specialization quadrant (represented up,
transfer down: \textbf{RQ1}). The KL-SFT arrows point up-and-left: it recovers transfer at a
represented-source cost (\textbf{RQ2}). Llama-Guard-3-1B sits at the origin because its cell is
\emph{degenerate, not robust}: its unmodified decision margin is a single constant on all $3{,}308$ rows
($-0.125$, sd $0$), so its macro-AP equals the source base rates by construction and no adaptation can
move it. We exclude it from the ensembling panel and treat it as uninformative rather than as evidence of
LoRA-resistance.}
\label{fig:adaptation-plane}
\end{figure}

\subsection{RQ1: ordinary SFT specializes released guards too (supported)}
\label{sec:adaptation-rq1}
\noindent
Across the panel, our SFT protocol raised represented-source macro-AP by \AdaHGain{} (equal-family mean;
LCB $\AdaHGainLCB>0$) while the gain was concentrated relative to held-out transfer (H\_conc \AdaHConc,
LCB $\AdaHConcLCB>0$), with the preregistered criterion $\text{LCB}(\text{H\_gain})>0$ and
$\text{LCB}(\text{H\_conc})>0$ both met. \Cref{tab:adaptation} and \Cref{fig:adaptation-plane} show the pattern is not a general-model
artifact: the released guards move the same way --- SFT buys represented ranking (e.g.\ ShieldGemma-2B
$+0.212$, Granite-Guard-2B $+0.139$) at a held-out cost (equal-family held-out change under SFT
\AdaHheldSFT). Fine-tuning a released guard on your data specializes it toward your sources, at a
bound-confirmed transfer cost, exactly as it does a general checkpoint.

\subsection{RQ2: KL-SFT retains transfer, but is \emph{not} a free improvement (not supported)}
\label{sec:adaptation-rq2}
\noindent
KL-SFT did preserve held-out transfer relative to SFT (H\_preserve \AdaHPreserve, LCB $\AdaHPreserveLCB>0$):
in \Cref{tab:adaptation} every KL-SFT held-out delta on the nine \emph{informative} checkpoints is less
negative (or more positive) than its SFT counterpart; on the degenerate Llama-Guard-3-1B cell both are
exactly zero. But the represented-source \emph{cost} of that retention, H\_cost \AdaHCost, has
LCB \AdaHCostLCB, which does not clear the preregistered non-inferiority margin of $-\AdaMargin$ --- so
the second criterion fails and RQ2 is \textbf{not supported}. KL-SFT trades represented-source adaptation
gain for transfer retention; on this panel it is a genuine trade, not a free lunch. This is a more
cautious verdict than the general-checkpoint KL control of \Cref{sec:actI-klsft} (where the represented
cost was smaller), and the confirmatory design reports it as such rather than reading the favorable
direction as a win.

\subsection{What this does and does not license}
\label{sec:adaptation-license}
\noindent
The within-checkpoint deltas are causal only for \emph{this recipe on this checkpoint}; the
general-vs-purpose contrast is a \emph{descriptive} blocked comparison (checkpoints are not randomly
assigned to a starting type), not a causal claim about starting type. The bounds are conditional on the
fixed model panel: the bootstrap resamples evaluation families and seeds but holds model identities
fixed, so with few families a single dominant family can move the equal-family mean; read H as a
fixed-panel summary, not a population estimate. One checkpoint, Llama-Guard-3-1B, shows essentially zero
movement under both SFT and KL-SFT (\Cref{tab:adaptation}): its $20$-token-pruned, embedding-tied output
head leaves the decision-token margins effectively unmoved by LoRA, so it contributes no adaptation
signal and should be read as a null/degenerate cell, not evidence of robustness. It is retained in the
equal-family means as a zero contribution; because a null family only pulls those means toward zero,
dropping it does not weaken either verdict (it would only raise the RQ1 gain/concentration bounds and
push the RQ2 cost bound further past the margin).

\begin{takeaway}{Adapting a released guard}
Fine-tuning a purpose-built guard on your data specializes it just like a general model (RQ1,
bound-confirmed): you gain represented ranking and lose transfer. KL-SFT buys the transfer back but
charges represented AP beyond the non-inferiority margin (RQ2 not supported) --- treat it as a
tradeoff dial, not a free upgrade, and re-measure both splits on your own data.
\end{takeaway}
